\input{estilo}
% Propuesta: el visibility timeout en el tiempo, con el caso de falla del
% texto. Arriba, el estado de m₁ en la cola; abajo, dos consumidores. a toma
% m₁ y falla sin borrarlo; mientras corre el timer, b pregunta y no lo ve;
% cuando el timer vence, m₁ vuelve a ser visible, b lo toma, lo procesa y lo
% borra. Todo sale del texto; el dibujo es nuevo.
\begin{document}
\begin{figura}\begin{tikzpicture}[
    tramo/.style={draw=tinta, minimum height=6mm, inner sep=0pt, outer sep=0pt,
      anchor=west, font=\footnotesize},
    proc/.style={line width=3pt, draw=suave!70},
  ]
  \def\yq{2.2} \def\ya{0.8} \def\yb{-0.6}
  % un tramo del estado de m₁: \tramo{x0}{x1}{relleno}{texto}
  \newcommand\tramo[4]{\node[tramo, fill=#3, minimum width={(#2-#1)*1cm}] at (#1,\yq) {#4};}
  % una llamada del consumidor a la cola y lo que vuelve: \llamada{x}{y}{pedido}{respuesta}
  \newcommand\llamada[4]{%
    \draw[msg] (#1-0.12,#2+0.12) -- node[etiq, left, pos=0.12] {#3} (#1-0.12,\yq-0.3);
    \draw[respuesta] (#1+0.12,\yq-0.3) -- node[etiq, right, pos=0.88] {#4} (#1+0.12,#2+0.12);}

  % los rótulos de las filas
  \node[anchor=east] at (-0.3,\yq) {\var{m}\textsubscript{1} en la cola};
  \node[anchor=east] at (-0.3,\ya) {consumidor \var{a}};
  \node[anchor=east] at (-0.3,\yb) {consumidor \var{b}};
  \foreach \y in {\ya, \yb} \draw[draw=suave, line width=0.4pt] (0,\y) -- (11.2,\y);

  % el estado de m₁
  \tramo{0}{1}{white}{visible}
  \tramo{1}{7}{suave!25}{invisible: corre el visibility timer, 60 \unid{s}}
  \tramo{7}{8.2}{white}{visible}
  \tramo{8.2}{10}{suave!25}{invisible}
  \node[anchor=west, font=\small] at (10.1,\yq) {borrado};
  \node[etiq, anchor=south] at (7,\yq+0.3) {vence el timer};
  \draw[marca] (7,\yq+0.3) -- (7,\yq-0.3);

  % a: lo toma, lo procesa y falla
  \llamada{1}{\ya}{receive}{\var{m}\textsubscript{1}}
  \draw[proc] (1.15,\ya) -- (3.6,\ya);
  \node[etiq, above] at (2.4,\ya) {procesa};
  \pic at (3.75,\ya) {rechazo};
  \node[etiq, text=acento, below=3pt] at (3.75,\ya) {falla, no borra};

  % b: no lo ve; después lo toma, lo procesa y lo borra
  \llamada{5.3}{\yb}{receive}{vacío}
  \llamada{8.2}{\yb}{receive}{\var{m}\textsubscript{1}}
  \draw[proc] (8.35,\yb) -- (9.85,\yb);
  \node[etiq, below] at (9.1,\yb) {procesa};
  \draw[msg] (10,\yb+0.12) -- node[etiq, right, pos=0.12] {delete} (10,\yq-0.3);

  \draw[->, draw=suave, line width=0.5pt] (0,-1.5) -- (11.2,-1.5) node[etiq, text=suave, below left] {tiempo};
\end{tikzpicture}\end{figura}
\end{document}
